## TEMP holding note — RQ1: does compartment count explain XGBoost's edge over the DNN? (2026-08-19)

**Which claim this checks**: the results chapter lists "the limited number of spatially
independent compartments" as one of several *possible, untested* reasons XGBoost outperforms the
DNN (`documentation/august_draft/5_Chapter_results_evaluation/g_written_draft_v2_redraft.tex`,
RQ1 spatial-block-performance section). If that mechanism is real, XGBoost's advantage over the
DNN should widen (not stay flat or shrink) when fewer independent compartments are available to
train on -- the same logic already used for the RQ3 47-compartment subsample check (rules out
compartment count alone as the explanation for 6survey's unresolvable Moran's $I$; see
`TEMP_rq2a_reduction_morans_i_2026-08-16.tex`), applied here to a different question.

**Status: complete (2026-08-19).**

### Method

**Code**: `models/baselines/rq1_compartment_subsample_check.py` (new file, this session). Run as:
```
PYTHONPATH=. .venv/bin/python -m models.baselines.rq1_compartment_subsample_check
```

**Pseudocode**:
```
load full 4survey table with Set3 (nested_set3_gated_terrain_wind_vif) terrain/wind features
merged in, via load_split_table_with_terrain() -- same loader run_dnn_env_terrain.py uses
  (this already applies the project's standard Age/yield-class/missing-feature filtering;
  232 compartments remain after filtering, not the raw 296)

for each condition in [full compartment set,
                       48-compartment subsample (seed 101),
                       48-compartment subsample (seed 102),
                       48-compartment subsample (seed 103)]:
    if subsampled:
        randomly choose 48 of the 232 compartments (numpy default_rng(seed), no replacement)
        filter the full table down to only rows in those compartments
    re-run spatial_block_split() fresh on whichever compartments are present (block_col='cpmt',
        60m buffer, split_seed=42 always -- so only WHICH compartments exist varies between
        conditions, not the block-shuffle logic itself)

    # both models see the exact same train/val/test rows within this condition
    fit DNN: standard fixed 128-128-128 architecture, default hyperparameters
        (lr=0.0001, batch_size=512, weight_decay=1e-5, seed=42), max_epochs=500,
        early_stopping_patience=40 -- identical to every other DNN fit in this project
    fit XGBoost: n_estimators=500, max_depth=6, learning_rate=0.02 -- the winning 4survey
        config from TEMP_rq1_xgb_hyperparameter_search_2026-08-16.tex, held fixed across every
        condition here (not re-tuned per condition, so only compartment count varies)

    record test R2 for both models and the gap (XGBoost R2 - DNN R2)

compare the full-compartment gap against the mean of the three 48-compartment gaps
```

**Design choices worth flagging**:
- Single `spatial_block` split per condition (not pooled 5-fold `spatial_block_kfold`) -- matches
  this project's established "coarse screen" convention for exploratory checks (e.g. the
  architecture sweep, `TEMP_rq1_architecture_sweep_results_2026-08-13.tex`), not the more
  expensive pooled-CV standard used for headline numbers. Read as directionally indicative, not a
  final-precision estimate.
- 48 compartments chosen to exactly match 6survey's real compartment count, so this result is
  directly comparable to 6survey's own real XGBoost-DNN gap (already in Table~\ref{tab:results-rq1}).
- 3 random seeds for which 48 compartments are kept (matching the RQ3 47-compartment precedent's
  "3 random seeds" convention), everything else (training seed, split-shuffle seed, both models'
  hyperparameters) held fixed, so compartment identity is the only thing varying between the three
  subsample replicates.
- XGBoost's config is NOT re-tuned per condition -- deliberate, isolates the effect of compartment
  count on a fixed model, rather than confounding it with "XGBoost gets retuned for less data too."

### Results

| Condition | n compartments | n train/val/test rows | DNN $R^2$ | XGBoost $R^2$ | Gap (XGB $-$ DNN) |
|---|---:|---|---:|---:|---:|
| Full compartment set | 232 | 129,536 / 47,056 / 46,956 | 0.6166 | 0.6402 | $+0.0236$ |
| Subsample, seed 101 | 48 | 22,224 / 9,400 / 11,796 | 0.6808 | 0.6370 | $-0.0438$ |
| Subsample, seed 102 | 48 | 24,692 / 9,212 / 9,004 | 0.5922 | 0.6209 | $+0.0287$ |
| Subsample, seed 103 | 48 | 30,684 / 10,840 / 10,716 | 0.6861 | 0.6014 | $-0.0847$ |

Subsampled (n=48) mean gap: $-0.0333$, SD $0.0469$ (across the 3 seeds), vs.\ the full-compartment
gap of $+0.0236$.

### Headline finding

**Not supported -- the opposite of what the hypothesis predicts.** If limited compartment count
were what favours XGBoost, the gap should widen (become more positive) as compartments shrink from
232 to 48. Instead the mean gap flips sign and goes negative ($+0.0236 \to -0.0333$), and the
individual seeds don't even agree on direction: XGBoost wins at seed 102 (+0.029, similar
magnitude to the full-compartment gap) but the DNN wins clearly at seeds 101 and 103 ($-0.044$,
$-0.085$). The seed-to-seed SD (0.047) is roughly double the full-compartment gap's own size --
at 48 compartments, which specific compartments happen to be sampled dominates the result far more
than any systematic model-family effect does.

**What this does and doesn't establish**: with only 3 replicates at n=48, this cannot rule out a
real but small effect hiding inside that noise -- it is not statistical proof of "no effect,
ever." What it does rule out is the specific, larger claim implied by citing this as a candidate
explanation for the real full-data XGBoost-DNN gap: there is no consistent, sign-stable widening
of XGBoost's advantage as compartment count drops, which is the direction the hypothesis actually
needs. On this evidence, compartment count is better read as adding instability to both models
alike at low N, not as a mechanism that systematically favours XGBoost. This mirrors the RQ3
47-compartment precedent's finding (`TEMP_rq2a_reduction_morans_i_2026-08-16.tex`): compartment
count alone does not explain effects seen on the real 48-compartment 6survey cohort -- something
about 6survey's actual composition (already documented in
`TEMP_rq1_cohort_composition_2026-08-15.tex`: lower, more homogeneous elevation range, a
genuinely different fitted CR curve) is the more likely explanation for the small residual gap
seen in that direction, not sample size in isolation.
